\section[title={2024-08-19 - Wohin geht diese
Generation?},reference={2024-08-19-Wohin geht diese Generation?}]

\startsectionlevel[title={Skeptische
Generation},reference={skeptische-generation}]

\startParagraph[reference=Vergleich,title=Vergleich]

\startitemize[packed]
\item
  1950er
\item
  Im sozialen Bewusstsein und Selbstbewusstsein kritischer, skeptischer,
  misstrauischer, illusionsloser als alle anderen Jugendgenerationen
  vorher.
\item
  Tolerant, ohne Pathos, Programme und Parolen.
\item
  Im privaten und sozialen Verhalten wirklichkeitsnäher,
  zugriffsbereiter und erfolgssicherer
\item
  Wird nie revolutionär reagieren -> trägt kein Bedürfnis elitäre
  Gemeinschaften zu stiften
\item
  Für Ältere scheint diese Gesellschaft als eine, die sich “totstellt“
  oder sich “tarnt“.
\item
  Sie setzt immer auf die Sicherheit.
\item
  Dienste der Generation der Gesamtgesellschaft und Öffentlichkeit
  liegen in den Tätigkeitsbereichen.
\item
  Technische Notwendigkeiten werden immer die höchste Überzeugungskraft
  haben.
\item
  wird alles Kollektive ablehnen
\item
  Geht aus der 68er Bewegung hervor
\stopitemize

\stopParagraph

\stopsectionlevel

\startsectionlevel[title={Kommende / Nachfolgende
Generation},reference={kommende-nachfolgende-generation}]

\startParagraph[reference=Vergleich,title=Vergleich]

\startitemize[packed]
\item
  Sezessionistisch (= trennt sich)
\item
  Ausbruch aus der Welt, Provokation
\item
  Proteste gegen \quotation{manipulierte} Freiheit u. Spontanität
\item
  Elitäre Reaktionen --> moralisch oder religiöse Rigorositäten
  (=Strenge/Härte)
\item
  Fantasie der \quotation{Ausbrüche} aus der \quotation{Wattewelt}
  Pädagogen u.ä. überlegen
\item
  Paragraph locked by Schüler Lana Lorbach
\stopitemize

\stopParagraph

\stopsectionlevel

\startsectionlevel[title={Aufgaben bis
Mittwoch},reference={aufgaben-bis-mittwoch}]

\startitemize[packed]
\item
  Markieren Si eim text, wie der Protagoniss\ldots{}
\stopitemize

\stopsectionlevel
